\label{sec:implem}
Beyond a mathematical framework that unlocks a more flexible use of orthogonal convolutions, we propose several design choices for an implementation that scales well to larger kernels, larger images and larger batch sizes. Although AOC includes the construction of BCOP and RKO kernels, our implementation improves the original ones at many stages. It results in an \textbf{8x} reduction of the original overhead in realistic settings. The implementation of AOC and several additional layers described in~\cref{sec:conv} is provided as an Open-source library called \fulllibname. 
%Our implementation will be integrated \libname. 
Other implementation details and empirical testing are discussed in \cref{ap:tech_details,ap:unit_testing}.

% \vspace{-3mm}

\vspace{-1mm}\paragraph{Fast implementation of the \Bconv.}\label{subsec:bconv_implem}

    To the best of our knowledge, the only differentiable implementation of the BCOP method is available in the reference code \cite{li_preventing_2019}. However, since there is no cuda kernel available for \Bconv, the authors relied on nested loops to perform all matrix multiplication to compute the resulting kernel. % Unfortunately, this approach prevents PyTorch from parallelizing the nested loops. %\sout{could be paralleled. Also, pytorch does not allow easily to parallelize cuda operators (that are already paralleled).}\franck{[ca parait pas hyper utile?]}
    % \paragraph{Hacking the 2D convolution to compute $\mathbconv$ operator.}
    We propose an efficient and parallelized implementation with a single operation. %However, our analysis have shown that this can be efficiently paralleled in a single operation. 
    %can be implemented as proposed in \cite{li_preventing_2019} by unrolling all the loops and explicit the full summation:
    % In order to efficiently implement this operation we can unroll all the loops and explicit the full summation:
    % \[
    %     R_{m,n,i,j} = \sum_{k \in 0}^{c-1}\sum_{i' \in 0}^{k-l-1}\sum_{j' \in 0}^{k-l-1} A_{m,k,i',j'} . B_{k,n,i-i',j-j'}
    % \]
    Inspired by \cite{wang_orthogonal_2020} and \cite{delattre2023efficient}, which aimed to compute $A\mathbconv A^T$ to prove orthogonality, we propose to replace the computation $B\mathbconv A$  by a convolution with zero padding between $B$ and $A^T$. This approach can also be seen as a specific case of convolutional einsum \cite{rabbani2024conv_einsum}.
    This operation can be rewritten by re-ordering the summation to use a 2D convolution at its core. The strategy is to use the 2D convolution to compute one output filter. Then, the batch dimension can be used to compute all output filters in parallel. The code is detailed in the \cref{fig:bconv_code}.

%             \begin{figure}
%                 \centering
%                     \begin{verbatim}
% def fast_matrix_conv(m1, m2, groups=1):
%     # m1 is m*n*k1*k2
%     # m2 is nb*m*l1*l2
%     m, n, k1, k2 = m1.shape
%     nb, mb, l1, l2 = m2.shape
%     assert m == mb * groups
%     # Rearrange m1 for conv
%     m1 = m1.transpose(0, 1)  # n*m*k1*k2
%     # Rearrange m2 for conv
%     m2 = m2.flip(-2, -1)
%     # Run convolution
%     # output shape nb*n*(k+l-1)*(k+l-1)
%     r2 = torch.nn.functional.conv2d(m1, m2, 
%            groups=groups, 
%            padding=(l1 - 1, l2 - 1))
%     # Rearrange result
%     # n*nb*(k+l-1)*(k+l-1)
%     return r2.transpose(0, 1)
%             \end{verbatim}
%                 \caption{Pytorch implementation of the algorithm. Tensorflow implementation differs only in the data format.}
%                 \label{fig:code}
%             \end{figure}

    
    % \begin{figure}
    %     \centering
    %     \includegraphics[width=\linewidth]{figures/ch3_aoc/benchmark/bloc_conv_speed_cin.png}
    %     \caption{kernel size is fixed while number of channels is increased. The runtime for 100 operation is reported on log scale. The previous implementation is very insensitive to the number of channels since the cost of the matrix multiplication is negligible against the cost of the convolution.}
    %     \label{fig:timeit_cin}
    % \end{figure}
    
    % \begin{figure}
    %     \centering
    %     \includegraphics[width=\linewidth]{figures/ch3_aoc/benchmark/bloc_conv_speed_k.png}
    %     \caption{number of channels is set to 16 whils kernel size in increased. The runtime for 100 operation is reported on log scale. This implementation scale better to larger kernel sizes.}
    %     \label{fig:timeit_cin}
    % \end{figure}
% \vspace{-3mm}

\vspace{-1mm}\paragraph{Reducing time complexity of BCOP.}
% \franck{tu ne parle pas de la paralleisation du calcul des projecteur symmetriques?}
    Beyond the efficient parallelism of the $\mathbconv$ operation, we propose to parallelize the whole kernel computation: the parametrization can be seen as the composition of many small kernels. Not only can those be created in parallel, but they can also be combined efficiently. As the $\mathbconv$ is an associative operation (\cref{prop:associativity}), we can leverage the parallel associative scan~\cite{zouzias2023parallel, daniel1986data} to parallelize the iterations of the original algorithm. The original $2*(k-s)$ sequential $\mathbconv$ operations can then be done in  $\mathcal{O} (\log (k-s)) $ iterations (\cref{fig:bcop_pscan}). This is unlocked in practice if $\mathbconv$ implementation supports batching. %Unfortunately, the batch dimension of our efficient implementation already uses the dimension originally dedicated to batching to compute the $c_o$ channels, so it cannot be used directly. 
    We propose to achieve this by using grouped \texttt{Conv2d} implementation: by concatenating $g$ kernels and setting $\text{groups}=g$, we can compute the batched $\mathbconv$ in parallel.
    %\franck{[ca veut dire quoi cette derniere phrase?]}However this implementation suffers the same limitation as the \textit{torch.bmm}: as it does not support heterogeneous kernel shapes.

% \vspace{-3mm}

\vspace{-1mm}\paragraph{Efficient implementation.}

    % \thib{faut faire revenir la def AOC $\kernel{K} = \kernel{K}_{\text{BCOP}} \mathbconv \kernel{K}_{\text{RKO}}$ du 3.2 on en a besoin ici}
    % By looking at \cref{def:method} one can notice that depending on values of $s$ and $k$ the parametrization can be simplified to either $\kernel{K} = \kernel{K}_{\text{BCOP}}$ or $\kernel{K} = \kernel{K}_{\text{RKO}}$. Furthermore we also proved that there exists cases where stride can be applied directly on a BCOP kernel \franck{cette phrase est disturbing, tu as affirmé depuis le debut que BCOP poucait pas gerer le stride, et tout a coup on dit le contraire....a reflechir}\thib{C'est le cas seulement quand ci < co et grace à la preuve \cref{proof:stride_opt}} without relying on the full parametrization (see \cref{proof:stride_opt}). The complete decision tree we use in our implementation can be found in the \cref{fig:method_simplify}. We proved the orthogonality of each branch through this paper.
    By examining \cref{def:AOC_kernel}, one observes that, depending on the values of $s$ and $k$, the parametrization can be simplified to either $ \kernel{K}_{AOC} = \kernel{K}_{\text{BCOP}} $ or  $\kernel{K}_{AOC} = \kernel{K}_{\text{RKO}}$. While BCOP is generally not suited for handling stride directly, we have identified specific cases -- namely when $c_i < c_o$ -- where stride can indeed be applied directly to a BCOP kernel (see \cref{proof:stride_opt}) without requiring the full parametrization. Although not proposed in \cite{li_preventing_2019}, this observation refines our overall characterization of BCOP’s limitations with stride, showing that exceptions exist under certain conditions. The complete decision tree used in our implementation is shown in \cref{fig:method_simplify}, with each branch’s orthogonality proved in \cref{sec:ap_proof_aoc_ortho,sec:ap_proof_rko,proof:stride_opt}.
    
    % \begin{figure}
    %     \centering
    %     \includegraphics[width=0.75\linewidth]{figures/ch3_aoc/illustrations/flowchart.drawio.png}
    %     \caption{\textbf{Our implementation automatically select the most efficient implementation}: under some circumstances, our method can be simplified by using only one of the two parametrizations initially required}
    %     \label{fig:flowchart}
    % \end{figure}
    
